\section{AIME I 1996}
        \begin{enumerate}
        	\item (\textit{AIME I 1996}) In a magic square, the sum of the three entries in any row, column, or diagonal is the same value. The figure shows four of the entries of a magic square. Find $x$. 


 
                    \begin{figure}[H]
                        \centering
                        \includegraphics[width=0.5\textwidth]{Problem/AIME_1996_I_1_Problem_Figure_1.png}
                    \end{figure}
                    
 \textbf{Answer: }200



\textbf{Solution}

Let's make a table.

 
 \[\begin{array}{|c|c|c|} \multicolumn{3}{c}{\text{Table}}\\\hline x&19&96\\\hline 1&a&b\\\hline c&d&e\\\hline \end{array}\]

 
 \begin{eqnarray*} x+19+96=x+1+c\Rightarrow c=19+96-1=114,\\ 114+96+a=x+1+114\Rightarrow a=x-95 \end{eqnarray*}

 
 \[\begin{array}{|c|c|c|} \multicolumn{3}{c}{\text{Table in progress}}\\\hline x&19&96\\\hline 1&x-95&b\\\hline 114&d&e\\\hline \end{array}\]

 
 \begin{eqnarray*}19+x-95+d=x+d-76=115+x\Rightarrow  d=191,\\ 114+191+e=x+115\Rightarrow e=x-190 \end{eqnarray*}
 \[\begin{array}{|c|c|c|} \multicolumn{3}{c}{\text{Table in progress}}\\\hline x&19&96\\\hline 1&x-95&b\\\hline 114&191&x-190\\\hline \end{array}\]

 
 \[3x-285=x+115\Rightarrow 2x=400\Rightarrow x=\boxed{200}\]

 





\textbf{Solution 2}

Use the table from above. Obviously $c = 114$. Hence $a+e = 115$. Similarly, $1+a = 96 + e \Rightarrow a = 95+e$.

 Substitute that into the first to get $2e = 20 \Rightarrow e=10$, so $a=105$, and so the value of $x$ is just $115+x = 210 + 105 \Rightarrow x = \boxed{200}$

 





\textbf{Solution 3}


 \[\begin{array}{|c|c|c|} \multicolumn{3}{c}{\text{Table}}\\\hline x&19&96\\\hline 1&a&b\\\hline c&d&e\\\hline \end{array}\]The formula 
 \[e=\frac{1+19}{2}\] can be used. Therefore, $e=10$. Similarly, 
 \[96=\frac{1+d}{2}\]So $d=191$.

 Now we have this table:
 \[\begin{array}{|c|c|c|} \multicolumn{3}{c}{\text{Table}}\\\hline x&19&96\\\hline 1&a&b\\\hline c&191&10\\\hline \end{array}\]By property of magic squares, observe that
 \[x+a+10=19+a+191\]The $a$'s cancel! We now have 
 \[x+10=19+191\]Thus $x=\boxed{200}.$

 




	\item (\textit{AIME I 1996}) For each real number $x$, let $\lfloor x \rfloor$ denote the greatest integer that does not exceed $x$. For how many positive integers $n$ is it true that $n<1000$ and that $\lfloor \log_{2} n \rfloor$ is a positive even integer?


 
 \textbf{Answer: }340



\textbf{Solution}

For integers $k$, we want $\lfloor \log_2 n\rfloor = 2k$, or $2k \le \log_2 n < 2k+1 \Longrightarrow 2^{2k} \le n < 2^{2k+1}$. Thus, $n$ must satisfy these \href{https://artofproblemsolving.com/wiki/index.php?title=Inequality}{inequalities} (since $n < 1000$):

 $4\leq n <8$
$16\leq n<32$
$64\leq n<128$


 $256\leq n<512$There are $4$ for the first inequality, $16$ for the second, $64$ for the third, and $256$ for the fourth, so the answer is $4+16+64+256=\boxed{340}$.

 




	\item (\textit{AIME I 1996}) Find the smallest positive integer $n$ for which the expansion of $(xy-3x+7y-21)^n$, after like terms have been collected, has at least 1996 terms.


 
 \textbf{Answer: }044



\textbf{Solution 1}

Using \href{https://artofproblemsolving.com/wiki/index.php?title=Simon\%27s\_Favorite\_Factoring\_Trick}{Simon's Favorite Factoring Trick}, we rewrite as $[(x+7)(y-3)]^n = (x+7)^n(y-3)^n$. Both \href{https://artofproblemsolving.com/wiki/index.php?title=Binomial\_expansion}{binomial expansions} will contain $n+1$ non-like terms; their product will contain $(n+1)^2$ terms, as each term will have an unique power of $x$ or $y$ and so none of the terms will need to be collected. Hence $(n+1)^2 \ge 1996$, the smallest square after $1996$ is $2025 = 45^2$, so our answer is $45 - 1 = \boxed{044}$.

 Alternatively, when $n = k$, the exponents of $x$ or $y$ in $x^i y^i$ can be any integer between $0$ and $k$ inclusive. Thus, when $n=1$, there are $(2)(2)$ terms and, when $n = k$, there are $(k+1)^2$ terms. Therefore, we need to find the smallest perfect square that is greater than $1996$. From trial and error, we get $44^2 = 1936$ and $45^2 = 2025$. Thus, $k = 45\rightarrow n = \boxed{044}$.

 





\textbf{Solution 2 (Generating Functions)}

Notice that the coefficients in the problem statement have no effect on how many unique terms there will be in the expansion. Therefore this problem is synonymous with finding the amount of terms in the expansion of $(xy+x+y+1)^n$ (we do this to simplify the problem).

 If we expand the exponent the expression becomes $\underbrace{(xy+x+y+1)\cdot(xy+x+y+1)\cdots(xy+x+y+1)}_{n\;times}$.

 This is equivalent to starting off with a $x^0y^0$ terms and choosing between $4$ options $n$ different times:

 $\bullet$ Adding nothing to either exponent (choosing $1$).

 $\bullet$ Adding $1$ to the $x$ exponent (choosing $x$).

 $\bullet$ Adding $1$ to the $y$ exponent (choosing $y$).

 $\bullet$ Adding $1$ to the $x$ exponent and adding $1$ to the $y$ exponent (choosing $xy$).

 Doing this $n$ times, you can see that you end up with a term in the form $k\cdot x^i\cdot y^j$ where $k$ is some coefficient (which we don't care about) and $0\le i,j\le n$ and $i,j \in \mathbb{N}$.

 Repeating this for all possible combinations of choices yields $n+1$ options for each of $i$ and $j$ which means there are a total of $(n+1)^2$ possible terms in the form $x^i\cdot y^j$. Therefore $(xy+x+y+1)^n$ has $(n+1)^2$ terms.

 $(n+1)^2 \ge 1996$ which yields $n=44$. $\boxed{044}$

 




	\item (\textit{AIME I 1996}) A wooden cube, whose edges are one centimeter long, rests on a horizontal surface. Illuminated by a point source of light that is $x$ centimeters directly above an upper vertex, the cube casts a shadow on the horizontal surface. The area of a shadow, which does not include the area beneath the cube is 48 square centimeters. Find the greatest integer that does not exceed $1000x$.


 
 \textbf{Answer: }166



\textbf{Solution}


                    \begin{figure}[H]
                        \centering
                        \includegraphics[width=0.5\textwidth]{Solution/AIME_1996_I_4_Solution_Figure_1.png}
                    \end{figure}
                    (Figure not to scale) The area of the square shadow base is $48 + 1 = 49$, and so the sides of the shadow are $7$. Using the similar triangles in blue, $\frac {x}{1} = \frac {1}{6}$, and $\left\lfloor 1000x \right\rfloor = \boxed{166}$.

 





\textbf{Solution (with more detail)}

Let a side of the shadow base not including the side length of cube be $y$. The triangle that $y$ is part of is similar to the triangle that $x$ is part of on the same side. 

 By similar triangles, $y = \frac {1}{x}$. The length of the complete side of the shadow base is $1+\frac {1}{x}$.

 The area of the shadow is $(1+\frac {1}{x})^2-1=48$.

 Solving this quadratic equation, continue as follows in Solution 1.

 
ANOTHER WAY is because we know 48 square cm is the area of the shadow, it must be the area of the larger square minus the area of the box, which is 1 square cm. So $(y+1)^2 - 1 = 48$ and $y = 6$.

 




	\item (\textit{AIME I 1996}) Suppose that the roots of $x^3+3x^2+4x-11=0$ are $a$, $b$, and $c$, and that the roots of $x^3+rx^2+sx+t=0$ are $a+b$, $b+c$, and $c+a$. Find $t$.


 
 \textbf{Answer: }023



\textbf{Solution 1}

By \href{https://artofproblemsolving.com/wiki/index.php?title=Vieta\%27s\_formulas}{Vieta's formulas} on the polynomial $P(x) = x^3+3x^2+4x-11 = (x-a)(x-b)(x-c) = 0$, we have $a + b + c = s = -3$, $ab + bc + ca = 4$, and $abc = 11$. Then

 $t = -(a+b)(b+c)(c+a) = -(s-a)(s-b)(s-c) = -(-3-a)(-3-b)(-3-c)$This is just the definition for $-P(-3) = \boxed{23}$.

 Alternatively, we can expand the expression to get
 \begin{align*} t &= -(-3-a)(-3-b)(-3-c)\\  &= (a+3)(b+3)(c+3)\\  &= abc + 3[ab + bc + ca] + 9[a + b + c] + 27\\ t &= 11 + 3(4) + 9(-3) + 27 = 23\end{align*}


 





\textbf{Solution  2}

Each term in the expansion of $(a+b)(b+c)(c+a)$ has a total degree of 3.  Another way to get terms with degree 3 is to multiply out $(a+b+c)(ab+bc+ca)$.  Expanding both of these expressions and comparing them shows that:

 $(a+b)(b+c)(c+a) = (ab+bc+ca)(a+b+c)-abc$$t = -(a+b)(b+c)(c+a) = abc-(ab+bc+ca)(a+b+c) = 11-(4)(-3) = 23$

 
A way to realize $(a+b)(b+c)(c+a) = (ab+bc+ca)(a+b+c)-abc$:

 $(a + b)(b+c)(c+a) = a^2b + a^2c + b^2a + b^2c + c^2a + c^2b + abc + abc$

 (Add an extra $abc$)

 $a^2b + a^2c + abc + b^2a + b^2c + abc + c^2a + c^2b + abc - abc =$

 $a(ab + bc + ac) + b(ab + bc + ac) + c(ab + bc + ac) - abc =$

 $(a + b + c)(ab + bc + ac) - abc = (-3)(4) - 11 = -23$.

 The value of $t$ is the negation of this, which is $-(-23) = \boxed{23}$

 





\textbf{Solution 3}

We have that $x^3+3x^2+4x-11=0$ for roots $a,b,c.$In the second cubic function $x^3+rx^2+sx+t=0,$ the roots are $a+b,b+c,c+a.$

 By Vieta's formulae, we see that $t= -(a+b)(b+c)(a+c).$As we know that the sum of the roots of the first polynomial, $a+b+c$ is $-3$ by applying Vieta's again.

 Using this fact, we can rewrite $t$ as $-(-3-a)(-3-b)(-3-c) = (a+3)(b+3)(c+3).$

 Seeing this, we can find the value of the product of the roots by applying this to the first equation.This can be done by setting $x'=a+3,$ so, from this, we see that we should plug in $x'-3$ for $x$ in $x^3+3x^2+4x-11=0$.

 After simplifying, we get that the polynomial is $x'^3 - 6x'^2 + 13x' - 23 = 0.$Given that the product of the roots of this equation is equivalent to our desired value for $t$, we can apply Vieta's formulae for a third time to find that $t = -(\frac{-23}{1}) = \boxed{23}.$

 




	\item (\textit{AIME I 1996}) In a five-team tournament, each team plays one game with every other team. Each team has a $50\%$ chance of winning any game it plays. (There are no ties.) Let $\dfrac{m}{n}$ be the probability that the tournament will produce neither an undefeated team nor a winless team, where $m$ and $n$ are relatively prime integers. Find $m+n$.


 
 \textbf{Answer: }049



\textbf{Solution 1}

We can use \href{https://artofproblemsolving.com/wiki/index.php?title=Complementary\_counting}{complementary counting}: finding the probability that at least one team wins all games or at least one team loses all games.

 No more than 1 team can win or lose all games, so at most one team can win all games and at most one team can lose all games.

 Now we use \href{https://artofproblemsolving.com/wiki/index.php?title=PIE}{PIE}:

 The probability that one team wins all games is $5\cdot \left(\frac{1}{2}\right)^4=\frac{5}{16}$.

 Similarity, the probability that one team loses all games is $\frac{5}{16}$.

 The probability that one team wins all games and another team loses all games is $\left(5\cdot \left(\frac{1}{2}\right)^4\right)\left(4\cdot \left(\frac{1}{2}\right)^3\right)=\frac{5}{32}$.

 $\frac{5}{16}+\frac{5}{16}-\frac{5}{32}=\frac{15}{32}$

 Since this is the opposite of the probability we want, we subtract that from 1 to get $\frac{17}{32}$.

 $17+32=\boxed{049}$

 





\textbf{Solution 2}

There are $\dbinom{5}{2} = 10$ games in total, and every game can either end in a win or a loss. Therefore, there are $2^{10} = 1024$ possible outcomes.

 Now, computing this probability directly seems a little hard, so let's compute the complement -- the probability that there is an undefeated team, a winless team, or both.

 The number of ways that we can have an undefeated team is $5 \cdot 2^6,$ as there are five ways to choose the team itself and $2^6$ outcomes for the games that are not concerning the undefeated team. Reversing all the wins to losses yields a winless team, so this situation is symmetrical to the first one. Therefore, there are $5 \cdot 2^6 + 5 \cdot 2^6$ ways for an undefeated or winless team.

 In the process, however, we have accidentally overcounted the scenario of both an undefeated and winless team. There are $5 \cdot 4 \cdot 2^3$ ways for this scenario to happen, because there are five ways to choose the undefeated team, four ways for the winless, and three games that don't concern either of the teams. Therefore, there are $5 \cdot 2^6 + 5 \cdot 2^6 - 5 \cdot 4 \cdot 2^3 = 480$ ways to have an undefeated and/or winless team, so there are $1024 - 480 = 544$ to not have any.

 Our probability is thus $\dfrac{544}{1024},$ which simplifies to $\dfrac{17}{32},$ for our answer of $17 + 32 = \boxed{049}.$

 Solution by Ilikeapos

 




	\item (\textit{AIME I 1996}) Two squares of a $7\times 7$ checkerboard are painted yellow, and the rest are painted green. Two color schemes are equivalent if one can be obtained from the other by applying a rotation in the plane board. How many inequivalent color schemes are possible?


 
 \textbf{Answer: }300



\textbf{Solution 1 (Generalized)}

There are ${49 \choose 2}$ possible ways to select two squares to be painted yellow. There are four possible ways to rotate each board. Given an arbitrary pair of yellow squares, these four rotations will either yield two or four equivalent but distinct boards. 

 <table><tbody><tr><td>          
                    \begin{figure}[H]
                        \centering
                        \includegraphics[width=0.5\textwidth]{Solution/AIME_1996_I_7_Solution_Figure_1.png}
                    \end{figure}
                              </td><td>
                    \begin{figure}[H]
                        \centering
                        \includegraphics[width=0.5\textwidth]{Solution/AIME_1996_I_7_Solution_Figure_1.png}
                    \end{figure}
                    </td></tr><tr><td><font style="font-size:85\%">For most pairs, there will be 
 three other equivalent boards.</font></td><td><font style="font-size:85\%">For those symmetric about the center, 
 there is only one other.</font></td></tr></tbody></table>Note that a pair of yellow squares will only yield $2$ distinct boards upon rotation \href{https://artofproblemsolving.com/wiki/index.php?title=Iff}{iff} the yellow squares are rotationally symmetric about the center square; there are $\frac{49-1}{2}=24$ such pairs. There are then ${49 \choose 2}-24$ pairs that yield $4$ distinct boards upon rotation; in other words, for each of the ${49 \choose 2}-24$ pairs, there are three other pairs that yield an equivalent board.

 Thus, the number of inequivalent boards is $\frac{{49 \choose 2} - 24}{4} + \frac{24}{2} = \boxed{300}$. For a $(2n+1) \times (2n+1)$ board, this argument generalizes to $n(n+1)(2n^2+2n+1)$ inequivalent configurations.

 





\textbf{Solution 2 (Casework)}

There are 4 cases: 1. The center square is occupied, in which there are $12$ cases.2. The center square isn't occupied and the two squares that are opposite to each other with respect to the center square, in which there are $12$ cases.3. The center square isn't occupied and the two squares can rotate to each other with a $90^{\circ}$ rotation with each other and with respect to the center square, in which case there are $12$ cases.4. And finally, the two squares can't rotate to each other with respect to the center square in which case there are $\dbinom{12}{2} \cdot \frac{16}{4} = 264$ cases. Add up all the values for each case to get $\boxed{300}$ as your answer.

 





\textbf{Solution 3 (Burnside’s Lemma)}

Consider the group $G=\mathbb{Z}/4\mathbb{Z}$ with addition acting on set $X$, where $X$ is the set of color schemes on the checkerboard. We now apply Burnside’s Lemma. Since $\mathbb{Z}/4\mathbb{Z}=\langle r \mid r^4=1\rangle$, we only need to consider 4 cases for $g$ an element of $\mathbb{Z}/4\mathbb{Z}$. 

 1. $g=1$. There are clearly ${49 \choose 2}$ cases.

 2. $g=r$. Since this is a $90$ degree rotation, there are no fixed points under this transformation. 

 3. $g=r^2$. There are $24$ cases (which can be manually checked).

 4. $g=r^3$. There are no cases since this is just Case $2$ but with a $-90$ degree rotation.

 Therefore, by Burnside’s Lemma there are $\frac{1176+24+0+0}{|G|}=\frac{1200}{4}=\boxed{300}$ inequivalent configurations (also known as orbits in group-theoretic terms).

 




	\item (\textit{AIME I 1996}) The harmonic mean of two positive integers is the reciprocal of the arithmetic mean of their reciprocals. For how many ordered pairs of positive integers $(x,y)$ with $x<y$ is the harmonic mean of $x$ and $y$ equal to $6^{20}$?


 
 \textbf{Answer: }799



\textbf{Solution}

The harmonic mean of $x$ and $y$ is equal to $\frac{1}{\frac{\frac{1}{x}+\frac{1}{y}}2} = \frac{2xy}{x+y}$, so we have $xy=(x+y)(3^{20}\cdot2^{19})$, and by \href{https://artofproblemsolving.com/wiki/index.php?title=SFFT}{SFFT}, $(x-3^{20}\cdot2^{19})(y-3^{20}\cdot2^{19})=3^{40}\cdot2^{38}$. Now, $3^{40}\cdot2^{38}$ has $41\cdot39=1599$ factors, one of which is the square root ($3^{20}2^{19}$). Since $x<y$, the answer is half of the remaining number of factors, which is $\frac{1599-1}{2}= \boxed{799}$.

 \begin{verbatim}
Note: Counting how many positive integer divisors of $n^2$ less than $n$ is also used in \href{https://artofproblemsolving.com/wiki/index.php?title=1995\_AIME\_Problems/Problem\_6}{1995 AIME Problems/Problem 6}.
\end{verbatim}




	\item (\textit{AIME I 1996}) A bored student walks down a hall that contains a row of closed lockers, numbered 1 to 1024. He opens the locker numbered 1, and then alternates between skipping and opening each locker thereafter. When he reaches the end of the hall, the student turns around and starts back. He opens the first closed locker he encounters, and then alternates between skipping and opening each closed locker thereafter. The student continues wandering back and forth in this manner until every locker is open. What is the number of the last locker he opens?


 
 \textbf{Answer: }342



\textbf{Solution 1}

On his first pass, he opens all of the odd lockers. So there are only even lockers closed. Then he opens the lockers that are multiples of $4$, leaving only lockers $2 \pmod{8}$ and $6 \pmod{8}$. Then he goes ahead and opens all lockers $2 \pmod {8}$, leaving lockers either $6 \pmod {16}$ or $14 \pmod {16}$. He then goes ahead and opens all lockers $14 \pmod {16}$, leaving the lockers either $6 \pmod {32}$ or $22 \pmod {32}$. He then goes ahead and opens all lockers $6 \pmod {32}$, leaving $22 \pmod {64}$ or $54 \pmod {64}$. He then opens $54 \pmod {64}$, leaving $22 \pmod {128}$ or $86 \pmod {128}$. He then opens $22 \pmod {128}$ and leaves $86 \pmod {256}$ and $214 \pmod {256}$. He then opens all $214 \pmod {256}$, so we have $86 \pmod {512}$ and $342 \pmod {512}$, leaving lockers $86, 342, 598$, and $854$, and he is at where he started again. He then opens $86$ and $598$, and then goes back and opens locker number $854$, leaving locker number $\boxed{342}$ untouched. He opens that locker.

 
Refer to the chart below for a more detailed view of the lockers that are left CLOSED on each turn, assuming that the lockers are numbered $1$ to $1024$ from left ($L$) to right ($R$) starting from $0 \pmod 2$. 

 \begin{verbatim}
Note: The general form in this chart is written differently, in the same way as Yiyj1 edited, to emphasize that the modulus is multiplied by two each time he makes a pass because he skips every other locker. 
\end{verbatim}
NOTICE THE PATTERNS: the least numbered and the greatest numbered lockers each change every two passes. The modulus changes every two passes as well. This speeds up calculation greatly.

 
For the first five passes:

 \begin{array}{lccccc}\text{Least number} & 2 & 2 & 6 & 6 & 22 \\\text{Greatest number} & 1024 & 1022 & 1022 & 1014 & 1014 \\\text{General form} & 0 \bmod 2 & -2 \bmod 4 & -2 \bmod 8 & -10 \bmod 16 & -10 \bmod 32 \\\text{Direction of next turn} & L & R & L & R & L \end{array}

 For the next five passes:

 \begin{array}{lccccc}\text{Least number} & 22 & 86 & 86 & 342 & 342 \\\text{Greatest number} & 982 & 982 & 854 & 854 \\\text{General form} & -42 \bmod 64 & -42 \bmod 128 & -170 \bmod 256 & -170 \bmod 512 \\\text{Direction of next turn} & R & L & R & L \end{array}

 





\textbf{Solution 2}

We can also solve this with recursion. Let $L_n$ be the last locker he opens given that he started with $2^n$ lockers. Let there be $2^n$ lockers. After he first reaches the end of the hallway, there are $2^{n-1}$ lockers remaining. There is a correspondence between these unopened lockers and if he began with $2^{n-1}$ lockers. The locker $y$ (if he started with $2^{n-1}$ lockers) corresponds to the locker $2^n+2-2y$ (if he started with $2^n$ lockers). It follows that $L_{n} = 2^{n} +2 -2L_{n-1}$ as they are corresponding lockers. We can compute $L_1=2$ and use the recursion to find $L_{10}=\boxed{342}$

 
I will illustrate this solution with a small example, and solving using a slightly different way taken from the first solution found on the Random Math Wiki page for this problem. Take for example if there were $2^2$ numbers: $[1 , 2 , 3 , 4]$. After the first pass, there are $2^2/2 = 2^1$ numbers left: $[2 , 4]$. Dealing with the new set on the second pass corresponds exactly to dealing with the set $[1 , 2]$ going from right to left ($[2 , 1]$), and we know that the last remaining locker of $[1 , 2]$ is $2$, which corresponds to $2$ in the original set. We know that $1$ in the bijection corresponds to $4$ in the original set after one pass, and $2$ (bijection) with $2$. In general terms, each number $n$ becomes $\frac{n}{2}$ because half the numbers are taken out in one pass, but you need to reverse the order of the new $2^{k-1}$ numbers so each $n$ becomes $2^{k-1} + 1 - \frac{n}{2}$ and because we consider the last number left after passing, $L_{k-1} = 2^{k-1} + 1 - \frac{L_k}{2}$ and because we want the original number that the bijection corresponds to, rearrange to get $L_{k} = 2^{k} +2 -2L_{k-1}$. We know the last number ($L_{k-1}$) of very small sequences (e.g. $2$), so for the doubled sequence the last number is the number that corresponds with the small sequence $2$, which is $2$. Then do recursion like the above solution states.

 





\textbf{Solution 3 (bash)}

List all the numbers from $1$ through $1024$, then do the process yourself!!! It will take about 25 minutes (if you don't start to see the pattern), but that's okay, eventually, you will get $\boxed{342}$.

 (Note: If you try to do this, first look through all the problems! -Guy)

 





\textbf{Solution 4 (bash but it takes less than 3 minutes)}

Note that whenever the student makes a pass starting on one end, the number at the other end of that line will remain. This is very important.

 Pass $0$:

 $1,2,3,4...,1024$

 Pass $1$:

 $2,4,6,...,1022, 1024$

 Pass $2$ (reversed, note that we take the second-to-last term of the previous sequence for every new sequence):

 $1022, 1018,...,(2+4), 2$

 Pass $3$:

 $6, (6+4+4),...,(1022-4-4), 1022$

 Pass $4$:

 $1014, (1014-16),..., (6+16), 6$

 Pass $5$:

 $22, (22+32),..., (1014-32), 1014$

 Pass $6$:

 $982, (982-64),..., (22 + 64), 22$

 Pass $7$:

 $86, (86+128),..., (982-128), 982$

 Pass $8$:

 $854, (854-256),..., (86+256), 86$

 Pass $9$:

 $342, (342+512)$

 Pass $10$:

 $342$

 The answer is therefore $\boxed{342}.$

 


 Edit from EthanSpoon, Doing this with only even numbers will make it faster.You'll see.02496

 Edit from Iamaimogoldmedalist. Doing some few rounds in your head till you get Multiples of 4 or possibly numbers of the form 8k+2 will enhance your speed.

 




	\item (\textit{AIME I 1996}) Find the smallest positive integer solution to $\tan{19x^{\circ}}=\dfrac{\cos{96^{\circ}}+\sin{96^{\circ}}}{\cos{96^{\circ}}-\sin{96^{\circ}}}$.


 
 \textbf{Answer: }159



\textbf{Solution}

$\dfrac{\cos{96^{\circ}}+\sin{96^{\circ}}}{\cos{96^{\circ}}-\sin{96^{\circ}}} =$ $\dfrac{\sin{186^{\circ}}+\sin{96^{\circ}}}{\sin{186^{\circ}}-\sin{96^{\circ}}} =$ $\dfrac{\sin{(141^{\circ}+45^{\circ})}+\sin{(141^{\circ}-45^{\circ})}}{\sin{(141^{\circ}+45^{\circ})}-\sin{(141^{\circ}-45^{\circ})}} =$ $\dfrac{2\sin{141^{\circ}}\cos{45^{\circ}}}{2\cos{141^{\circ}}\sin{45^{\circ}}} = \tan{141^{\circ}}$. 

 The period of the \href{https://artofproblemsolving.com/wiki/index.php?title=Trigonometry\#Tangent}{tangent} function is $180^\circ$, and the tangent function is \href{https://artofproblemsolving.com/wiki/index.php?title=One-to-one}{one-to-one} over each period of its domain. 

 Thus, $19x \equiv 141 \pmod{180}$. 

 Since $19^2 \equiv 361 \equiv 1 \pmod{180}$, multiplying both sides by $19$ yields $x \equiv 141 \cdot 19 \equiv (140+1)(18+1) \equiv 0+140+18+1 \equiv 159 \pmod{180}$. 

 Therefore, the smallest positive solution is $x = \boxed{159}$.

 





\textbf{Solution 2}

$\dfrac{\cos{96^{\circ}}+\sin{96^{\circ}}}{\cos{96^{\circ}}-\sin{96^{\circ}}} = \dfrac{1 + \tan{96^{\circ}}}{1-\tan{96^{\circ}}}$which is the same as $\dfrac{\tan{45^{\circ}} + \tan{96^{\circ}}}{1-\tan{45^{\circ}}\tan{96^{\circ}}} = \tan{141{^\circ}}$.

 So $19x = 141 +180n$, for some integer $n$.Multiplying by $19$ gives $x \equiv 141 \cdot 19 \equiv 2679 \equiv 159 \pmod{180}$.The smallest positive solution of this is $x = \boxed{159}$

 





\textbf{Solution 3 (Only sine and cosine sum formulas)}

It seems reasonable to assume that $\dfrac{\cos{96^{\circ}}+\sin{96^{\circ}}}{\cos{96^{\circ}}-\sin{96^{\circ}}} = \tan{\theta}$ for some angle $\theta$. This means 
 \[\dfrac{\alpha (\cos{96^{\circ}}+\sin{96^{\circ}})}{\alpha (\cos{96^{\circ}}-\sin{96^{\circ}})} = \frac{\sin{\theta}}{\cos{\theta}}\] for some constant $\alpha$. We can set $\alpha (\cos{96^{\circ}}+\sin{96^{\circ}}) = \sin{\theta}$.Note that if we have $\alpha$ equal to both the sine and cosine of an angle, we can use the sine sum formula (and the cosine sum formula on the denominator). So, since $\sin{45^{\circ}} = \cos{45^{\circ}} = \frac{\sqrt{2}}{2}$, if $\alpha = \frac{\sqrt{2}}{2}$ we have 
 \[\alpha (\cos{96^{\circ}} + \sin{96^{\circ}}) = \cos{96^{\circ}} \frac{\sqrt{2}}{2} + \sin{96^{\circ}} \frac{\sqrt{2}}{2} = \cos{96^{\circ}} \sin{45^{\circ}} + \sin{96^{\circ}} \cos{45^{\circ}} = \sin({45^{\circ} + 96^{\circ}}) = \sin{141^{\circ}}\]from the sine sum formula. For the denominator, from the cosine sum formula, we have 
 \[\alpha (\cos{96^{\circ}} - \sin{96^{\circ}}) = \cos{96^{\circ}} \frac{\sqrt{2}}{2} + \sin{96^{\circ}} \frac{\sqrt{2}}{2} = \cos{96^{\circ}} \cos{45^{\circ}} + \sin{96^{\circ}} \sin{45^{\circ}} = \cos({96^{\circ}  + 45^{\circ}}) = \cos{141^{\circ}}.\]This means $\theta = 141^{\circ},$ so $19x = 141 + 180k$ for some positive integer $k$ (since the period of tangent is $180^{\circ}$), or $19 x \equiv 141 \pmod{180}$. Note that the inverse of $19$ modulo $180$ is itself as $19^2 \equiv 361 \equiv 1 \pmod {180}$, so multiplying this congruence by $19$ on both sides gives $x \equiv 2679 \equiv 159 \pmod{180}.$ For the smallest possible $x$, we take $x = \boxed{159}.$

 





\textbf{Solution 4}

Multiplying the numerator and denominator of the right-hand side by$\cos(96^{\circ})+\sin(96^{\circ})$, we get 

 ${\tan(19x^{\circ})}  ={\frac{\cos(96^{\circ}) +\sin(96^{\circ})}{\cos(96^{\circ})-\sin(96^{\circ})}}\times{\frac{\cos(96^{\circ})+\sin(96^{\circ})}{\cos(96^{\circ})+\sin(96^{\circ})}} \\  ={\frac{(\cos(96^{\circ})+\sin(96^{\circ}))^2}{\cos^2(96^{\circ})-\sin^2(96^{\circ})}} \\  ={\frac{\cos^2(96^{\circ}) + 2\cos(96^{\circ})\sin(96^{\circ}) + \sin^2(96^{\circ})}{\cos(192^{\circ})}} \\  ={\frac{1+\sin(192^{\circ})}{\cos(192^{\circ})}}$

 Using the fact that $\tan(\theta)=\frac{\sin(\theta)}{\cos(\theta)}$, we get $\tan(19x^{\circ})=\frac{\sin(19x^{\circ})}{\cos(19x^{\circ})}=\frac{1+\sin(192^{\circ})}{\cos(192^{\circ})}$.

 Cross-multiplying, we find that$\sin(19x^{\circ})\cos(192^{\circ})=\cos(19x^{\circ})+\cos(19x^{\circ})\sin(192^{\circ})$.

 Rearranging the equation gives us$\cos(19x^{\circ})=\sin(19x^{\circ})\cos(192^{\circ})-\cos(19x^{\circ})\sin(192^{\circ})$which leads us to $\cos(19x^{\circ})=\sin(19x-192^{\circ})$ by the sine difference formula.

 Using the identity that $\cos(\theta)=\sin(90^{\circ}-\theta)$, we find that $\sin(90-19x^{\circ})=\sin(19x-192^{\circ})$. 

 Therefore, $90-19x \equiv 19x-192 \pmod{360}$, or $38x \equiv 282 \pmod{360}$. 

 We know that $38 \times 9=342$ and $38 \times 10 \equiv 20 \pmod{360}$ (by simple arithmetic). To "make" $282$ we subtract $10$ three times from $9$, giving us $-21$. 

 Finally, because $360|38 \times 180$, we can add $180$ to get that$x=180-21=\boxed{159}$ which is the final answer.

 


 





\textbf{Solution 5}

Notice that $\sin 45^\circ = \cos 45^\circ$. Therefore, all terms in the numerator and the denominator could be multiplied by either $\sin 45^\circ$ or $\cos 45^\circ$.
 \begin{align*}     \frac{\cos 96^\circ + \sin 96^\circ}{\cos 96^\circ - \sin 96^\circ}     &= \frac{\cos 96^\circ \sin 45^\circ + \sin 96^\circ \cos 45^\circ}{\cos 96^\circ \cos 45^\circ - \sin 96^\circ \sin 45^\circ} \\[0.5em]     &= \frac{\sin(96 + 45)^\circ}{\cos(96 + 45)^\circ} \\[0.5em]     &= \tan 141^\circ \end{align*}
Therefore, $\tan19x^\circ = \tan 141^\circ$. In other words, $19x = 141 + 180k$ for some integer $k$.
 \begin{align*}     141 + 180k &\equiv 0 \pmod{19} \\     8 + 9k &\equiv 0 \pmod{19} \\     9k &\equiv 11 \pmod{19} \\     9k &\equiv 11, 30, 49, 68, 87, 106, 125, 144 \pmod{19} \\ \end{align*}
Substituting $k$, the smallest value of $19x$ is $141 + 180 \cdot 16$, which is $\boxed{159}$.

 




	\item (\textit{AIME I 1996}) Let $\mathrm {P}$ be the product of the roots of $z^6+z^4+z^3+z^2+1=0$ that have a positive imaginary part, and suppose that $\mathrm {P}=r(\cos{\theta^{\circ}}+i\sin{\theta^{\circ}})$, where $0<r$ and $0\leq \theta <360$. Find $\theta$.


 
 \textbf{Answer: }276



\textbf{Solution 1}


 \begin{eqnarray*} 0 &=& z^6 - z + z^4 + z^3 + z^2 + z + 1 = z(z^5 - 1) + \frac{z^5-1}{z-1}\\ 0 &=& \frac{(z^5 - 1)(z(z-1)+1)}{z-1} = \frac{(z^2-z+1)(z^5-1)}{z-1} \end{eqnarray*}

 
Thus $z^5 = 1, z \neq 1 \Longrightarrow z = \mathrm{cis}\ 72, 144, 216, 288$,

 or $z^2 - z + 1 = 0 \Longrightarrow z = \frac{1 \pm \sqrt{-3}}{2} = \mathrm{cis}\ 60, 300$ 

 (see \href{https://artofproblemsolving.com/wiki/index.php?title=Cis}{cis}).

 Discarding the roots with negative imaginary parts (leaving us with $\mathrm{cis} \theta,\ 0 < \theta < 180$), we are left with $\mathrm{cis}\ 60, 72, 144$; their product is $P = \mathrm{cis} (60 + 72 + 144) = \mathrm{cis} \boxed{276}$.

 \begin{verbatim}
To avoid using the quadratic formula,
\end{verbatim}$0 = \frac{(z^5 - 1)(z(z-1)+1)}{z-1} = \frac{(z^2-z+1)(z^5-1)}{z-1} = \frac{(z^3+1)(z^5-1)}{(z+1)(z-1)}$

 Thus $z^5 = 1, z \neq 1, -1 \Longrightarrow z = \mathrm{cis}\ 72, 144, 216, 288$,

 or $z^3 = -1, z \neq -1 \Longrightarrow z = \mathrm{cis}\ 60, 300$ 

 (see \href{https://artofproblemsolving.com/wiki/index.php?title=Roots\_of\_unity}{Roots of unity} for how to solve these equations)

 Then proceed as above.

 





\textbf{Solution 2}

Let $w =$ the fifth \href{https://artofproblemsolving.com/wiki/index.php?title=Roots\_of\_unity}{roots of unity}, except for $1$. Then $w^6 + w^4 + w^3 + w^2 + 1 = w^4 + w^3 + w^2 + w + 1 = 0$, and since both sides have the fifth roots of unity as roots, we have $z^4 + z^3 + z^2 + z + 1  | z^6 + z^4 + z^3 + z^2 + 1$. Long division quickly gives the other factor to be $z^2 - z + 1$. The solution follows as above.

 





\textbf{Solution 3}

Divide through by $z^3$. We get the equation $z^3 + \frac {1}{z^3} + z + \frac {1}{z} + 1 = 0$. Let $x = z + \frac {1}{z}$. Then $z^3 + \frac {1}{z^3} = x^3 - 3x$. Our equation is then $x^3 - 3x + x + 1 = x^3 - 2x + 1 = (x - 1)(x^2 + x - 1) = 0$, with solutions $x = 1, \frac { - 1\pm\sqrt {5}}{2}$. For $x = 1$, we get $z = \text{cis}60,\text{cis}300$. For $x = \frac { - 1 + \sqrt {5}}{2}$, we get $z = \text{cis}{72},\text{cis}{288}$ (using exponential form of $\cos$). For $x = \frac { - 1 - \sqrt {5}}{2}$, we get $z = \text{cis}144,\text{cis}216$. The ones with positive imaginary parts are ones where $0\le\theta\le180$, so we have $60 + 72 + 144 = \boxed{276}$.

 





\textbf{Solution 4 (Fakesolve)}

This is just a slight variation of Solution 1.

 We start off by adding $z^5$ to both sides, to get a neat geometric sequence with $a = 1$ and $r = z$, which gives us $\frac{z^7 - 1}{z - 1} = z^5$. From here, multiply by $z - 1$ to both sides, noting that then $z \neq \cos 0 + i\sin 0$ since, then we are multiplying by $0$ which makes it undefined. We now note that $z^7 - 1 = z^6 - z^5 \implies z^7 - z^6 + z^5 = 1$. (This is the part that it becomes almost identical to Solution 1). Factor $z^5$ from the LHS, to get $z^5( z^2 - z + 1) = 1$. Call the set of roots from $z^5$ as $A$, and set of roots from $z^2 - z + 1$ as $B$. We are to find $|A \cup B| = A + B - |A \cap B|$. (Essentially, we are to find the roots that are not common to both equations/sets, or else we are overcounting a root two times, rather than once. Try out some equation to see where this might apply) Thankfully in this case, none of the roots are overcounted. From here, proceed with Solution 1.

 





\textbf{Solution 5}

We recognize that $z^6+z^5+z^4+z^3+z^2+z+1=\frac{z^7-1}{z-1}$ and strongly wish for the equation to turn into so. Thankfully, we can do this by simultaneously adding and subtracting $z^5$ and $z$ as shown below.$z^6+z^5+z^4+z^3+z^2+z+1-(z^5+z)=0\implies \frac{z^7-1}{z-1}-(z^5+z)=0$

 Now, knowing that $z=1$ is not a root, we multiply by $z-1$ to obtain$z^7-1-(z-1)(z^5+z)=0\implies z^7-1-(z^6+z^2-z^5-z)=0\implies z^7-z^6+z^5-z^2+z-1=0$Now, we see the $z^2+z-1$ and it reminds us of the sum of two cubes. Cleverly factoring, we obtain that..$z^5(z^2-z)+z^5-(z^2-z+1)=0\implies z^5(z^2-z+1)-(z^2-z+1)=0\implies (z^5-1)(z^2-z+1)=0$. 

 Now, it is clear that we have two cases to consider.

 Case $1$: $z^5-1=0$We obtain that $z^5=1$ or $z^5=e^{2\pi{n}{i}}$Obviously, the answers to this case are $e^{ia}, a\in{\frac{2\pi}{5}, \frac{4\pi}{5}}$

 Case $2$: $z^2-z+1=0$Completing the square and then algebra allows us to find that $z=\frac{1}{2}\pm\frac{i\sqrt{3}}{2}$ which has $\arg$ $\frac{\pi}{3}$

 Hence, the answer is $\frac{18\pi}{15}+\frac{5\pi}{15}=\frac{23\pi}{15}\cdot\frac{180}{\pi}=\boxed{276}$

 





\textbf{Solution 6}

Add 1 to both sides of the equation to get $x^6+x^4+x^3+x^2+1+1=1$. We can rearrange to find that $(x^6+x^3+1)+(x^4+x^2+1)=1$. Then, using sum of a geometric series, $\frac{x^9-1}{x^3-1}+\frac{x^6-1}{x^2-1}=1$.

 Combining the two terms of the LHS, we get that $\frac{x^{11}-x^9-x^2+1+x^9-x^6-x^3+1}{x^5-x^3-x^2+1}=1$, so $x^{11}-x^6-x^3-x^2+2=x^5-x^3-x^2+1$, and simplifying, we see that $x^{11}-x^6-x^5+1=0$, so by SFFT, $(x^6-1)(x^5-1)=0$. Then, the roots of our polynomial are the fifth and sixth roots of unity. However, looking back at our expression when it had fractions, we realize that if $x$ is a second or third root of unity, it would cause a denominator to be zero, so the roots of our polynomial are the fifth and sixth roots of unity that are not the second or third roots of unity. The only roots in this category are $\mathrm{cis} (60, 72, 144)$, so our desired sum is $\boxed{276}$, and we are done.

 -coolak

 





\textbf{Solution 7}

Take accout of the $Z^3$ and use the formula $(x^n-1)(x^m+x^{m-n}+...+1) = x^{m+n}-1$

 Then factorise the function to solve $Z^5=1$, $Z^3 \neq 1$, $Z^2 \neq 1$

 




	\item (\textit{AIME I 1996}) For each permutation $a_1,a_2,a_3,\cdots,a_{10}$ of the integers $1,2,3,\cdots,10$, form the sum


 $|a_1-a_2|+|a_3-a_4|+|a_5-a_6|+|a_7-a_8|+|a_9-a_{10}|$.


 The average value of all such sums can be written in the form $\dfrac{p}{q}$, where $p$ and $q$ are relatively prime positive integers. Find $p+q$.


 
 \textbf{Answer: }058



\textbf{Solution 1}

Because of symmetry, we may find all the possible values for $|a_n - a_{n - 1}|$ and multiply by the number of times this value appears. Each occurs $5 \cdot 8!$, because if you fix $a_n$ and $a_{n + 1}$ there are still $8!$ spots for the others and you can do this $5$ times because there are $5$ places $a_n$ and $a_{n + 1}$ can be.

 To find all possible values for $|a_n - a_{n - 1}|$ we have to compute
 \begin{eqnarray*} |1 - 10| + |1 - 9| + \ldots + |1 - 2|\\ + |2 - 10| + \ldots + |2 - 3| + |2 - 1|\\ + \ldots\\ + |10 - 9| \end{eqnarray*}

 This is equivalent to

 
 \[2\sum\limits_{k = 1}^{9}\sum\limits_{j = 1}^{k}j = 330\]

 The total number of permutations is $10!$, so the average value is $\frac {330 \cdot  8!  \cdot  5}{10!} = \frac {55}{3}$, and $m+n = \boxed{058}$.

 Note from SuperJJ: Another way to think about this is that the total sum is 330. Because you have a total of $10 \cdot 9$ possible "pairs" or "differences", the average difference is $\frac{330}{90} = \frac{11}{3}$. Because the question asks for the expected value of the sum of $5$ of these pairs / differences, our answer is just $\frac{11}{3} \cdot 5 = {\frac{55}{3}}$.

 





\textbf{Solution 2}

\href{https://artofproblemsolving.com/wiki/index.php?title=Without\_loss\_of\_generality}{Without loss of generality}, let $a_1 > a_2;\, a_3 > a_4;\, \ldots ;\, a_9 > a_{10}$. We may do this because all sums obtained from these paired sequences are also obtained in another $2^5-1$ ways by permuting the adjacent terms $\{a_1,a_2\},\{a_3,a_4\}, \cdots , \{a_9, a_{10}\}$, and thus are canceled when the average is taken. 

 So now we only have to form the sum $S= (a_1 + a_3 + a_5 + a_7 + a_9) - (a_2 + a_4 + a_6 + a_8 + a_{10})$. Due to the symmetry of this situation, we only need to compute the \href{https://artofproblemsolving.com/wiki/index.php?title=Expected\_value}{expected value} of the result. $10$ must always be the greatest number in its pair; $9$ will be the greater number in its pair $\frac{8}{9}$of the time and the lesser number $\frac 19$ of the time; $8$ will be the greater number in its pair $\frac 79$ of the time and the lesser $\frac 29$ of the time; and so forth. Each number either adds or subtracts from the sum depending upon whether it is one of the five greater or five lesser numbers in the pairs, respectively. Thus

 
 \begin{align*} \overline{S} &= 10 + \left(\frac{8}{9} \cdot 9\right) - \left(\frac{1}{9} \cdot 9\right) + \left(\frac{7}{9} \cdot 8\right) - \left(\frac 29 \cdot 8\right) + \cdots + \left(\frac{1}{9} \cdot 2\right) - \left(\frac{8}{9} \cdot 2\right) - 1 \\ &= \frac{9 \cdot 10 + 7 \cdot 9 + 5 \cdot 8 + 3 \cdot 7 + 1 \cdot 6 - 1 \cdot 5 - 3 \cdot 4 - 5 \cdot 3 - 7 \cdot 2 - 9 \cdot 1}{9} \\ &= \frac{55}{3} \end{align*}


 And the answer is $m+n = \boxed{058}$.

 





\textbf{Solution 3}

Similar to Solution 1, we can find the average value of $|a_2 - a_1|$, and multiply this by 5 due to symmetry.  And again due to symmetry, we can arbitrarily choose $a_2 > a_1$.  Thus there are $\binom{10}{2} = 45$ ways to pick the two values of $a_2$ and $a_1$ from the set $\{1, 2, 3, 4, 5, 6, 7, 8, 9, 10\}$ such that $a_2 > a_1$.  First fix $a_2 = 10$, and vary $a_1$ from $1$ to $9$.  Then fix $a_2 = 9$, and vary $a_1$ from $1$ to $8$.  Continue, and you find that the sum of these $45$ ways to pick $|a_2 - a_1|$ is:

 $\sum\limits_{k = 1}^{9}\sum\limits_{j = 1}^{k}j = 45+36+28+21+15+10+6+3+1 = 165$.  

 Thus, each term contributes on average $\frac{165}{45}$, and the sum will be five times this, or $\frac{165}{9} = \frac{55}{3}$.

 The final answer is $p+q = \boxed{058}$.

 





\textbf{Solution 4(Expected Value)}

We use expected value of one of the sums, say $|a_2 - a_1|$, since all five are similar, so the average or expected value of the sum is just 5 times the expected value of one of them.

 To pick two random expected numbers from 1 to 10, we use a well known expected value trick by tying the two ends 1 and 10 with an extra number 0.* This gives 11 spaces, and we distribute 3 gaps evenly to choose our two numbers. We get $\frac{11}{3}$ and $\frac{22}{3}$ giving an absolute difference of $\frac{11}{3}$. Therefore, the average expected value of the sum of the five would be $\frac{55}{3}$. This gives $55+3=\boxed{058}$.

 
                \begin{itemize}
                    \item This is like a bendable stick that has equally spaced marks "end-1-2-3-4-5-6-7-8-9-10-end", and the ends are bent together giving a circle with 11 equally spaced marks(the ends produce 1 mark together).
                \end{itemize}


                





\textbf{Solution 4a (Different Take on Same Method, Maybe Easier to Understand?)}

You can also think of this as a straight number line starting at 0 and ending at 11. Pick two integers from 1 to 10, which means there will be a gap from 0 to the smaller number, the smaller number to the larger number, and the larger number to 11. The expected length of the gaps are all the same, so the length of each gap is $\frac{11}{3}$ (total distance from 0 to 11 divided by 3 gaps).

 \begin{verbatim}
Why start from 0 and end at 11? This is to preserve symmetry. 
\end{verbatim}\begin{verbatim}
Let the number line start at 1 and end at 10. If the smaller number was 1 and the larger number was 10, making the greatest possible middle gap, the leftmost gap would be 0 and the rightmost gap would be 0, too. However, the smallest gap possible between the smaller number and larger number is 1 (eg. 1 and 2, 2 and 3, ...)
\end{verbatim}\begin{verbatim}
Every gap must have a minimum size of 1, because if you don't use 0 and 11, the outer gaps are allowed to be 0, but the middle gap is forced to be at least 1. Because the middle gap has a strict head start, it is fundamentally different from the others. You can no longer divide the total evenly because the rules aren't the same. By choosing 0 and 11, every single gap now has a minimum size of exactly 1. They now share identical constraints. You're still only picking integers from 1 to 10, but the way you count the differences is even now.
\end{verbatim}Continue as in Solution 4 to find the final answer.

 




	\item (\textit{AIME I 1996}) In triangle $ABC$, $AB=\sqrt{30}$, $AC=\sqrt{6}$, and $BC=\sqrt{15}$. There is a point $D$ for which $\overline{AD}$ bisects $\overline{BC}$, and $\angle ADB$ is a right angle. The ratio


 
 \[\dfrac{\text{Area}(\triangle ADB)}{\text{Area}(\triangle ABC)}\]


 can be written in the form $\dfrac{m}{n}$, where $m$ and $n$ are relatively prime positive integers. Find $m+n$.


 
 \textbf{Answer: }065



\textbf{Solution 1}

Proceed by Barycentrics. Let $D = (t : 1 : 1)$, so we want $\left|\frac{1}{t+2}\right|$. Let $E$ be the foot of the altitude from $A$. Note $E = (0 : -\frac{9}{2} : \frac{39}{2})$. $ADBE$ is cyclic, so $D, E$ are on circle 
 \[-a^2yz-b^2xz-c^2xy + (x+y+z)(wz) = 0\] Substituting $E$, we find $w = -\frac{9}{2}$. Thus, substituting in the coordinates for $D$, we find 
 \[-\frac{9}{2}(t+2) = a^2+t(b^2+c^2) = 15+36t \implies t = -\frac{16}{27} \implies \frac{1}{t+2} = \frac{27}{38} \implies \boxed{065}\]

 





\textbf{Solution 2}


                    \begin{figure}[H]
                        \centering
                        \includegraphics[width=0.5\textwidth]{Solution/AIME_1996_I_13_Solution_Figure_1.png}
                    \end{figure}
                    Let $E$ be the midpoint of $\overline{BC}$. Since $BE = EC$, then $\triangle ABE$ and $\triangle AEC$ share the same height and have equal bases, and thus have the same area. Similarly, $\triangle BDE$ and $BAE$ share the same height, and have bases in the ratio $DE : AE$, so $\frac{[BDE]}{[BAE]} = \frac{DE}{AE}$ (see \href{https://artofproblemsolving.com/wiki/index.php?title=Area\_ratios}{area ratios}). Now,

 $\dfrac{[ADB]}{[ABC]} = \frac{[ABE] + [BDE]}{2[ABE]} = \frac{1}{2} + \frac{DE}{2AE}.$By \href{https://artofproblemsolving.com/wiki/index.php?title=Stewart\%27s\_Theorem}{Stewart's Theorem}, $AE = \frac{\sqrt{2(AB^2 + AC^2) - BC^2}}2 = \frac{\sqrt {57}}{2}$, and by the \href{https://artofproblemsolving.com/wiki/index.php?title=Pythagorean\_Theorem}{Pythagorean Theorem} on $\triangle ABD, \triangle EBD$,

 
 \begin{align*} BD^2 + \left(DE + \frac {\sqrt{57}}2\right)^2 &= 30 \\ BD^2 + DE^2 &= \frac{15}{4} \\ \end{align*}


 Subtracting the two equations yields $DE\sqrt{57} + \frac{57}{4} = \frac{105}{4} \Longrightarrow DE = \frac{12}{\sqrt{57}}$. Then $\frac mn = \frac{1}{2} + \frac{DE}{2AE} = \frac{1}{2} + \frac{\frac{12}{\sqrt{57}}}{2 \cdot \frac{\sqrt{57}}{2}} = \frac{27}{38}$, and $m+n = \boxed{065}$.

 


 





\textbf{Solution 3}

Because the problem asks for a ratio, we can divide each side length by $\sqrt{3}$ to make things simpler. We now have a triangle with sides $\sqrt{10}$, $\sqrt{5}$, and $\sqrt{2}$.

 We use the same graph as above.

 Draw perpendicular from $C$ to $AE$. Denote this point as $F$. We know that $DE = EF = x$ and $BD = CF = z$ and also let $AE = y$.

 Using Pythagorean theorem, we get three equations,

 <center>$(y+x)^2 + z^2 = 10$$(y-x)^2 + z^2 = 2$

 $x^2 + z^2 = \frac{5}{4}$</center>Adding the first and second, we obtain $x^2 + y^2 + z^2 = 6$, and then subtracting the third from this we find that $y = \frac{\sqrt{19}}{2}$. (Note, we could have used \href{https://artofproblemsolving.com/wiki/index.php?title=Stewart\%27s\_Theorem}{Stewart's Theorem} to achieve this result). 

 Subtracting the first and second, we see that $xy = 2$, and then we find that $x = \frac{4}{\sqrt{19}}$ 

 Using base ratios, we then quickly find that the desired ratio is $\frac{27}{38}$ so our answer is $\boxed{065}$

 





\textbf{Solution 4 (Trig Bash)}


                    \begin{figure}[H]
                        \centering
                        \includegraphics[width=0.5\textwidth]{Solution/AIME_1996_I_13_Solution_Figure_1.png}
                    \end{figure}
                    First, extend side \(AC\) to a point \(F\) such that \(AF \perp BF\).

 We begin by using the Law of Cosines to find \(\cos \angle ACB\):
 \[(\sqrt{6})^2 + (\sqrt{15})^2 - 2(\sqrt{6})(\sqrt{15})\cos \angle ACB = (\sqrt{30})^2.\]Simplifying,
 \[6 + 15 - 2\sqrt{90}\cos \angle ACB = 30,\]which gives
 \[-6\sqrt{10}\cos \angle ACB = 9 \quad \Rightarrow \quad \cos \angle ACB = -\frac{3}{2\sqrt{10}}.\]

 Using the Pythagorean identity,
 \[\sin \angle ACB = \sqrt{1 - \cos^2 \angle ACB} = \frac{\sqrt{310}}{20}.\]

 Let \([ABC]\) denote the area of \(\triangle ABC\). Using the area formula \(\frac{1}{2}ab\sin C\),
 \[[ABC] = \frac{1}{2}(\sqrt{6})(\sqrt{15})\left(\frac{\sqrt{310}}{20}\right) = \frac{3\sqrt{31}}{4}.\]

 Next, we use the Law of Cosines to find \(AE\):
 \[AE^2 = \left(\frac{\sqrt{15}}{2}\right)^2 + (\sqrt{6})^2 - 2\left(\frac{\sqrt{15}}{2}\right)(\sqrt{6})\left(-\frac{3}{2\sqrt{10}}\right).\]Simplifying,
 \[AE^2 = \frac{39}{4} + \frac{9}{2} = \frac{57}{4},\]so
 \[AE = \frac{\sqrt{57}}{2}.\]

 Since \(E\) is the midpoint of \(BC\), we have \([ACE] = [ABE] = \frac{3\sqrt{31}}{8}\).

 Solving for \(\sin \angle BEA\),
 \[\frac{1}{2}\left(\frac{\sqrt{15}}{2}\right)\left(\frac{\sqrt{57}}{2}\right)\sin \angle BEA = \frac{3\sqrt{31}}{8},\]which simplifies to
 \[\sqrt{95}\sin \angle BEA = \sqrt{31},\]so
 \[\sin \angle BEA = \sqrt{\frac{31}{95}}.\]

 Using a Pythagorean identity,
 \[\cos \angle BEA = \frac{8}{\sqrt{95}}.\]

 Now, solving for \(DE\),
 \[\frac{DE}{\frac{\sqrt{15}}{2}} = \frac{8}{\sqrt{95}} \quad \Rightarrow \quad DE = \frac{4\sqrt{15}}{\sqrt{95}} = 4\sqrt{\frac{3}{19}}.\]

 We now find the area of \(\triangle BDE\):
 \[[BDE] = \frac{1}{2}\left(4\sqrt{\frac{3}{19}}\right)\left(\frac{\sqrt{15}}{2}\right)\left(\sqrt{\frac{31}{95}}\right) = \frac{6\sqrt{31}}{38}.\]

 Adding areas,
 \[[ADB] = [ABE] + [BDE] = \frac{3\sqrt{31}}{8} + \frac{6\sqrt{31}}{38} = \frac{81\sqrt{31}}{152}.\]

 Taking the ratio,
 \[\frac{[ADB]}{[ABC]} = \frac{\frac{81\sqrt{31}}{152}}{\frac{3\sqrt{31}}{4}} = \frac{27}{38}.\]

 Since the problem asks for \(m+n\), we have \(27 + 38 = \boxed{065}.\)

 


 




	\item (\textit{AIME I 1996}) A $150\times 324\times 375$ rectangular solid is made by gluing together $1\times 1\times 1$ cubes. An internal diagonal of this solid passes through the interiors of how many of the $1\times 1\times 1$ cubes?


 
 \textbf{Answer: }768



\textbf{Solution 1}

Place one corner of the solid at $(0,0,0)$ and let $(a,b,c)$ be the general coordinates of the diagonally opposite corner of the rectangle, where $a, b, c \in \mathbb{Z_{+}}$.

 We consider the \href{https://artofproblemsolving.com/wiki/index.php?title=Vector}{vector} equation of the diagonal segment represented by:$(x, y, z) =m (a, b, c)$, where $m \in \mathbb{R}, 0 < m \le 1$.

 Consider a point on the diagonal with coordinates $(ma, mb, mc)$. We have 3 key observations as this point moves from $(0,0,0)$ towards $(a,b,c)$:

 
            \begin{enumerate}
                \item When exactly one of $ma$, $mb$, $mc$ becomes a positive integer, the point crosses one of the \href{https://artofproblemsolving.com/wiki/index.php?title=Face}{faces} of a $1 \times 1 \times 1$ cube from the outside to the inside of the cube.
            \end{enumerate}
            
            \begin{enumerate}
                \item When exactly two of $ma$, $mb$, $mc$ become positive integers, the point enters a new cube through one of the edges of the cube from the outside to the inside of the cube.
            \end{enumerate}
            
            \begin{enumerate}
                \item When all three $ma$, $mb$, $mc$ become positive integers, the point enters a cube through one of its vertices from the outside to the inside of the cube.
            \end{enumerate}
            
The number of cubes the diagonal passes is equal to the number of points on the diagonal that has one or more positive integers as coordinates.

 If we slice the solid up by the $x$-planes defined by $x=1,2,3,4, \ldots, a$, the diagonal will cut these $x$-planes exactly $a$ times (plane of $x=0$ is not considered since $m \ne 0$). Similar arguments for slices along $y$-planes and $z$-planes give diagonal cuts of $b$, and $c$ times respectively. The total cuts by the diagonal is therefore $a+b+c$, if we can ensure that no more than $1$ positive integer is present in the x, y, or z coordinate in all points $(ma,mb,mc)$ on the diagonal. Note that point $(a,b,c)$ is already one such exception.

 But for each diagonal point $(ma,mb,mc)$ with 2 (or more) \href{https://artofproblemsolving.com/wiki/index.php?title=Positive}{positive} \href{https://artofproblemsolving.com/wiki/index.php?title=Integer}{integers} occurring at the same time, $a+b+c$ counts the number of cube passes as $2$ instead of $1$ for each such point.  There are $\gcd(a,b)+\gcd(b,c)+\gcd(c,a)$ points in such over-counting. We therefore subtract one time such over-counting from $a+b+c$.

 And for each diagonal point $(ma,mb,mc)$ with exactly $3$ integers occurring at the same time, $a+b+c$ counts the number of cube passes as $3$ instead of $1$; ie $a+b+c$ over-counts each of such points by $2$. But since $\gcd(a,b)+\gcd(b,c)+\gcd(c,a)$ already subtracted three times for the case of $3$ integers occurring at the same time (since there are $3$ of these gcd terms that represent all combinations of 3 edges of a cube meeting at a vertex), we  have the final count for each such point as $1 = 3-3+k \Rightarrow k=1$, where $k$ is our correction term. That is, we need to add $k=1$ time $\gcd(a,b,c)$ back to account for the case of 3 simultaneous integers.

 Therefore, the total diagonal cube passes is:$D = a+b+c-\left[ \gcd(a,b)+\gcd(b,c)+\gcd(c,a) \right]+\gcd(a,b,c)$.

 For $(a,b,c) = (150, 324, 375)$, we have:$\gcd(150,324)=6$, $\gcd(324,375)=3$, $\gcd(150,375)=75$, $\gcd(150,324,375)=3$.

 Therefore$D = 150+324+375-(6+3+75)+3 = \boxed{768}$.

 





\textbf{Solution 2}

Consider a point travelling across the internal diagonal, and let the internal diagonal have a length of $d$. The point enters a new unit cube in the $x,y,z$ dimensions at multiples of $\frac{d}{150}, \frac{d}{324}, \frac{d}{375}$ respectively. We proceed by using PIE.

 The point enters a new cube in the $x$ dimension $150$ times, in the $y$ dimension $324$ times and in the $z$ dimension, $375$ times.

 The point enters a new cube in the $x$ and $y$ dimensions whenever a multiple of $\frac{d}{150}$ equals a multiple of $\frac{d}{324}$. This occurs $\gcd(150, 324)$ times. Similarly, a point enters a new cube in the $y,z$ dimensions $\gcd(324, 375)$ times and a point enters a new cube in the $z,x$ dimensions $\gcd(375, 150)$ times.

 The point enters a new cube in the $x,y$ and $z$ dimensions whenever some multiples of $\frac{d}{150}, \frac{d}{324}, \frac{d}{375}$ are equal. This occurs $\gcd(150, 324, 375)$ times.

 The total number of unit cubes entered is then $150+324+375-[\gcd(150, 324)+\gcd(324, 375)+\gcd(375, 150)] + \gcd(150, 324, 375) = \boxed{768}$

 




	\item (\textit{AIME I 1996}) In parallelogram $ABCD,$ let $O$ be the intersection of diagonals $\overline{AC}$ and $\overline{BD}$. Angles $CAB$ and $DBC$ are each twice as large as angle $DBA,$ and angle $ACB$ is $r$ times as large as angle $AOB$. Find the greatest integer that does not exceed $1000r$.


 
 \textbf{Answer: }777



\textbf{Solution 1 (trigonometry)}


                    \begin{figure}[H]
                        \centering
                        \includegraphics[width=0.5\textwidth]{Solution/AIME_1996_I_15_Solution_Figure_1.png}
                    \end{figure}
                    Let $\theta = \angle DBA$. Then $\angle CAB = \angle DBC = 2\theta$, $\angle AOB = 180 - 3\theta$, and $\angle ACB = 180 - 5\theta$. Since $ABCD$ is a parallelogram, it follows that $OA = OC$. By the \href{https://artofproblemsolving.com/wiki/index.php?title=Law\_of\_Sines}{Law of Sines} on $\triangle ABO,\, \triangle BCO$,

 $\frac{\sin \angle CBO}{OC} = \frac{\sin \angle ACB}{OB} \quad \text{and} \quad \frac{\sin \angle DBA}{OC} = \frac{\sin \angle BAC}{OB}.$Dividing the two equalities yields 

 
 \[\frac{\sin 2\theta}{\sin \theta} = \frac{\sin (180 - 5\theta)}{\sin 2\theta} \Longrightarrow \sin^2 2\theta = \sin 5\theta \sin \theta.\]

 Pythagorean and product-to-sum identities yield

 
 \[1 - \cos^2 2 \theta = \frac{\cos 4\theta - \cos 6 \theta}{2},\]

 and the double and triple angle ($\cos 3x = 4\cos^3 x - 3\cos x$) formulas further simplify this to

 
 \[4\cos^3 2\theta - 4\cos^2 2\theta - 3\cos 2\theta + 3 = (4\cos^2 2\theta - 3)(\cos 2\theta - 1) = 0\] 

 The only value of $\theta$ that fits in this context comes from $4 \cos^2 2\theta - 3 = 0 \Longrightarrow \cos 2\theta = \frac{\sqrt{3}}{2} \Longrightarrow \theta = 15^{\circ}$. The answer is $\lfloor 1000r \rfloor = \left\lfloor 1000 \cdot \frac{180 - 5\theta}{180 - 3\theta} \right\rfloor = \left \lfloor \frac{7000}{9} \right \rfloor = \boxed{777}$.

 





\textbf{Solution 2 (trigonometry)}

Define $\theta$ as above. Since $\angle CAB = \angle CBO$, it follows that $\triangle COB \sim \triangle CBA$, and so $\frac{CO}{BC} = \frac{BC}{AC} \Longrightarrow BC^2 = AC \cdot CO = 2CO^2 \Longrightarrow BC = CO\sqrt{2}$. The \href{https://artofproblemsolving.com/wiki/index.php?title=Law\_of\_Sines}{Law of Sines} on $\triangle BOC$ yields that 

 
 \[\frac{BC}{CO} = \frac{\sin (180-3\theta)}{\sin 2\theta} = \frac{\sin 3\theta}{\sin 2\theta} = \sqrt{2}\]

 Expanding using the sine double and triple angle formulas, we have 

 
 \[2\sqrt {2} \sin \theta \cos \theta = \sin\theta( - 4\sin ^2 \theta + 3) \Longrightarrow \sin \theta\left(4\cos^2\theta - 2\sqrt {2} \cos \theta - 1\right) = 0.\]

 By the \href{https://artofproblemsolving.com/wiki/index.php?title=Quadratic\_formula}{quadratic formula}, we have $\cos \theta = \frac {2\sqrt {2} \pm \sqrt {8 + 4 \cdot 1 \cdot 4}}{8} = \frac {\sqrt {6} \pm \sqrt {2}}{4}$, so $\theta = 15^{\circ}$ (as the other roots are too large to make sense in context). The answer follows as above.

 





\textbf{Solution 3}

We will focus on $\triangle ABC$. Let $\angle ABO = x$, so $\angle BAO = \angle OBC = 2x$. Draw the \href{https://artofproblemsolving.com/wiki/index.php?title=Perpendicular}{perpendicular} from $C$ intersecting $AB$ at $H$. \href{https://artofproblemsolving.com/wiki/index.php?title=Without\_loss\_of\_generality}{Without loss of generality}, let $AO = CO = 1$. Then $HO = 1$, since $O$ is the \href{https://artofproblemsolving.com/wiki/index.php?title=Circumcenter}{circumcenter} of $\triangle AHC$. Then $\angle OHA = 2x$. 

 By the Exterior Angle Theorem, $\angle COB = 3x$ and $\angle COH = 4x$. That implies that $\angle HOB = x$. That makes $HO = HB = 1$. Then since by AA ($\angle HBC = \angle HOC = 3x$ and reflexive on $\angle OCB$), $\triangle OCB \sim \triangle BCA$.

 $\frac {CO}{BC} = \frac {BC}{AC} \implies 2 = BC^2 = \implies BC = \sqrt {2}.$Then by the \href{https://artofproblemsolving.com/wiki/index.php?title=Pythagorean\_Theorem}{Pythagorean Theorem}, $1^2 + HC^2 = \left(\sqrt {2}\right)^2\implies HC = 1$. That makes $\triangle HOC$ equilateral. Then $\angle HOC = 4x = 60 \implies x = 15$. The answer follows as above. 

 




\end{enumerate}